\normalsize \begin{longtable} {|p{0.06\textwidth}|p{0.25\textwidth}|p{0.7\textwidth}|} \caption{DM Science Roles for Rubin Observatory Data Management Operations \label{tab:dmsciroles}}\\ 
\hline 
\textbf{WBS}&\textbf{Role Title}&\textbf{Role Description} \\ \hline
{3.13}&{}&{Support the community's use of the survey data products, science pipelines, and data access services (e.g., Rubin Science Platform). Ingest feedback needed to ensure scientific productivity.} \\ \hline
{3.13a}&{Lead Community Scientist}&{Serves as strategic and tactical lead of the Community Science team, and the primary interface between the user community  and the operations team (e.g. by facilitating the Users Committee). Oversees user training and support, documentation, implementation of the processes used for evaluating applications for computing time and special survey cadences.  Serves as responsible lead of Science Communication, with sign-off on data product documentation readiness (recommending to the AD of System Performance when the data product documentation is ready for a release).  Reports to the AD of System Performance. This is a relatively senior role to be filled by a scientist with the respect and recognition of the community. } \\ \hline
{3.13b.6}&{Deputy Lead Community Scientist}&{Assists with leadership of the CST and assumes the lead role when needed. Maintains a high level of expertise with CST policies, workflows, activities, and interfaces. As a Community Scientist, supports the broad international scientific community of users through various communication channels (e.g., online Q\&A forums, workshops, tutorials), and collaborates closely with staff in Observatory Operations and Data Production to diagnose and respond to issues identified by the scientific community.  Has scientific experience and expertise in at least one of (or spanning) the Rubin Observatory science pillars. } \\ \hline
{3.13b}&{Community Scientist - NOIRLab Astronomer}&{Supports the broad international scientific community of users through various communication channels (e.g., online Q\&A forums, workshops, tutorials). Collaborates closely with staff in Observatory Operations and Data Production to diagnose and respond to issues identified by the scientific community.  These generalists should have significant depth of experience and expertise spanning the Rubin Observatory science pillars (and especially the transient universe, solar system, and Milky Way). The Astronomer role is filled by a scientist on the NOIRLab Astronomer track, and includes 50\% science time. } \\ \hline
{3.13b.1}&{Community Scientist Engagement PD- NOIRLab}&{Supports the broad international scientific community of users through various communication channels (e.g., online Q\&A forums, workshops, tutorials). Collaborates closely with staff in Observatory Operations and Data Production to diagnose and respond to issues identified by the scientific community.  These generalists should have experience and expertise spanning the Rubin Observatory science pillars (and especially the transient universe, solar system, and Milky Way). This role is to be filled by a postdoc and comes with 50\% science time.} \\ \hline
{3.13b.2}&{Community Scientist - NOIRLab Scientist}&{Supports the broad international scientific community of users through various communication channels (e.g., online Q\&A forums, workshops, tutorials). Collaborates closely with staff in Observatory Operations and Data Production to diagnose and respond to issues identified by the scientific community.  These generalists should have experience and expertise spanning the Rubin Observatory science pillars (and especially the transient universe, solar system, and Milky Way). } \\ \hline
{3.13b.3}&{Community Scientist for Chile - NOIRLab Scientist}&{Supports the Chilean scientific community of users through various communication channels (e.g., online Q\&A forums, workshops, tutorials). Collaborates closely with staff in Observatory Operations and Data Production to diagnose and respond to issues identified by Chilean scientific community.  These generalists should have experience and expertise spanning the Rubin Observatory science pillars (and especially the transient universe and Milky Way, the two most active science pillars in Chile). } \\ \hline
{3.13b.4}&{Surge Community Scientist analysis priority II-1 - NOIRLab Scientist}&{Supports the Chilean scientific community of users through various communication channels (e.g., online Q\&A forums, workshops, tutorials). Collaborates closely with staff in Observatory Operations and Data Production to diagnose and respond to issues identified by Chilean scientific community.  These generalists should have experience and expertise spanning the Rubin Observatory science pillars (and especially the transient universe and Milky Way, the two most active science pillars in Chile). } \\ \hline
{3.13b.5}&{Surge Community Scientist analysis priority II-1 - NOIRLab Scientist}&{Supports the Chilean scientific community of users through various communication channels (e.g., online Q\&A forums, workshops, tutorials). Collaborates closely with staff in Observatory Operations and Data Production to diagnose and respond to issues identified by Chilean scientific community.  These generalists should have experience and expertise spanning the Rubin Observatory science pillars (and especially the transient universe and Milky Way, the two most active science pillars in Chile). } \\ \hline
{3.13c.1}&{Community Scientist - SLAC Ops}&{Supports the broad international scientific community of users through various communication channels (e.g., online Q\&A forums, workshops, tutorials). Collaborates closely with staff in Observatory Operations and Data Production to diagnose and respond to issues identified by the scientific community.  These generalists should have experience and expertise spanning the Rubin Observatory science pillars (and especially the study of dark matter and dark energy).} \\ \hline
{3.13c.1}&{Community Scientist - SLAC Ops}&{Supports the broad international scientific community of users through various communication channels (e.g., online Q\&A forums, workshops, tutorials). Collaborates closely with staff in Observatory Operations and Data Production to diagnose and respond to issues identified by the scientific community.  These generalists should have experience and expertise spanning the Rubin Observatory science pillars (and especially the study of dark matter and dark energy).} \\ \hline
{3.13c.2}&{Surge Community Scientist - SLAC}&{Additional effort to help prepare for DP1 and DP2, cover potential pre-ops staff turnover. 50\% science time to cement connection to DESC and enable recruitment.} \\ \hline
{3.13c.3}&{Community Scientist - SLAC Res}&{Supports the broad international scientific community of users through various communication channels (e.g., online Q\&A forums, workshops, tutorials). Collaborates closely with staff in Observatory Operations and Data Production to diagnose and respond to issues identified by the scientific community.  These generalists should have experience and expertise spanning the Rubin Observatory science pillars (and especially the study of dark matter and dark energy).} \\ \hline
{3.13c.4}&{Community Scientist - Fermilab Ops}&{Supports the broad international scientific community of users through various communication channels (e.g., online Q\&A forums, workshops, tutorials). Collaborates closely with staff in Observatory Operations and Data Production to diagnose and respond to issues identified by the scientific community.  These generalists should have experience and expertise spanning the Rubin Observatory science pillars (and especially the study of dark matter and dark energy).} \\ \hline
{3.13c.5}&{Community Scientist - Fermilab Res}&{Supports the broad international scientific community of users through various communication channels (e.g., online Q\&A forums, workshops, tutorials). Collaborates closely with staff in Observatory Operations and Data Production to diagnose and respond to issues identified by the scientific community.  These generalists should have experience and expertise spanning the Rubin Observatory science pillars (and especially the study of dark matter and dark energy).} \\ \hline
{3.13d}&{Community Scientist - Secondments}&{Supports the broad international scientific community of users through various communication channels (e.g., online Q\&A forums, workshops, tutorials). Collaborates closely with staff in Observatory Operations and Data Production to diagnose and respond to issues identified by the scientific community.  These scientists should have experience and expertise that collectively span the Rubin Observatory science pillars, and can be postdocs circulating through NOIRLab and SLAC on secondments. Staffing profile gives an approximate target team size.} \\ \hline
{3.13e.1}&{Community Scientist (Citizen Science) - SLAC}&{Responsible for ensuring that citizen science programs enhance the scientific performance of the Rubin Observatory. Provides support for members of the science community to set up, execute, and retrieve results from citizen science programs that are an integral part of their Rubin-based research. Evaluates and logs citizen science inquiries and Rubin Observatory-Zooniverse integration troubleshooting. Participates in the Zooniverse framework technology refresh cycle that typically occurs every 3-5 years. Within the Rubin Citizen Science Working Group,  works with the EPO Scientist, Citizen Science Developer and Citizen Science Advisor to guide the development of citizen science programs. Leverages other Community Scientists and Documentation staff to help with the curation of User Generated data products (including their SDQA) so as to render them suitable for EPO programs at the EDC. Responsible for quality assurance of the EPO Data Products pipeline into the EDC and the scientific validity of EPO deliverables. Works with the EPO Scientist to select and develop citizen science projects that require customizations not provided via the Zooniverse Project Builder Tool.} \\ \hline
{3.13e.2}&{Community Scientist (Citizen Science) - SLAC Surge}&{Responsible for ensuring that citizen science programs enhance the scientific performance of the Rubin Observatory. Provides support for members of the science community to set up, execute, and retrieve results from citizen science programs that are an integral part of their Rubin-based research. Evaluates and logs citizen science inquiries and Rubin Observatory-Zooniverse integration troubleshooting. Participates in the Zooniverse framework technology refresh cycle that typically occurs every 3-5 years. Within the Rubin Citizen Science Working Group,  works with the EPO Scientist, Citizen Science Developer and Citizen Science Advisor to guide the development of citizen science programs. Leverages other Community Scientists and Documentation staff to help with the curation of User Generated data products (including their SDQA) so as to render them suitable for EPO programs at the EDC. Responsible for quality assurance of the EPO Data Products pipeline into the EDC and the scientific validity of EPO deliverables. Works with the EPO Scientist to select and develop citizen science projects that require customizations not provided via the Zooniverse Project Builder Tool.} \\ \hline
{3.13f}&{Community Documentation - NOIRLab}&{Coordinates and participates in the delivery of user-facing scientific documentation related to the Rubin Observatory data products, and software, as well as scientific information on the Rubin Observatory website.  Oversees and makes decisions on user-facing documentation tools and methods, with approval from the Lead Community Scientist. Reports to the Community Scientist.} \\ \hline
{3.13g}&{Community Documentation - SLAC}&{Coordinates and participates in the delivery of user-facing scientific documentation related to the Rubin Observatory data products, and software, as well as scientific information on the Rubin Observatory website.  Oversees and makes decisions on user-facing documentation tools and methods, with approval from the Lead Community Scientist. Reports to the Community Scientist.} \\ \hline
{3.13g}&{Community Documentation - SLAC}&{Coordinates and participates in the delivery of user-facing scientific documentation related to the Rubin Observatory data products, and software, as well as scientific information on the Rubin Observatory website.  Oversees and makes decisions on user-facing documentation tools and methods, with approval from the Lead Community Scientist. Reports to the Community Scientist.} \\ \hline
\end{longtable} \normalsize
